\documentclass[11pt,a4paper]{article}
\usepackage{amsmath,amssymb,amsthm}
\usepackage[margin=2.5cm]{geometry}
\usepackage{hyperref}

\newtheorem{definition}{Definition}[section]
\newtheorem{theorem}[definition]{Theorem}
\newtheorem{proposition}[definition]{Proposition}
\newtheorem{lemma}[definition]{Lemma}
\newtheorem{corollary}[definition]{Corollary}
\newtheorem{conjecture}[definition]{Conjecture}
\newtheorem{remark}[definition]{Remark}

\title{Hyperseed Formalization:\\
  The Leaky Transcension Hypothesis}
\author{ProtomegaTron (automated formalization)}
\date{2026-09-18}

\begin{document}
\maketitle

\begin{abstract}
We formalize the intellectual structure of Ben Goertzel's Substack article
``The Leaky Transcension Hypothesis'' (2026-04-25) within the Hyperseed
ontology. The article extends John Smart's Transcension Hypothesis with a
``leaky'' variant: advanced civilizations transcend inward but leak
unavoidably at the boundaries. Key mappings: transcension as fiber
deepening, leakage as fiber-to-base projection, boundary effects as
fiber-base interface, Fermi paradox as fiber invisibility, efficiency
as fiber density, AGI transcension as fiber-deepening architecture,
and leakage management as projection control.
\end{abstract}

\tableofcontents

%% =================================================================
\section{Transcension as fiber deepening}
\label{sec:deepening}
%% =================================================================

\begin{theorem}[Inward convergence --- claims 1--4, 18]
\label{thm:deepening}
Smart's Transcension Hypothesis: advanced civilizations converge toward
inner space (computation, simulation, black-hole computing) rather than
outer space (Dyson spheres, interstellar colonization). Technology
consistently moves toward miniaturization, densification, and efficiency.

In Hyperseed terms, transcension is \emph{fiber deepening}: moving to
higher fiber depth (richer internal structure) rather than wider base
space (more spatial extent):
\[
  \text{Transcension:} \quad \dim(F) \uparrow \quad \text{while} \quad
  \text{vol}(B) \text{ constant or } \downarrow
\]

The fiber gets deeper --- more computational structure, richer internal
representation, denser information processing --- while the base footprint
(physical extent in space) stays constant or shrinks.

The efficiency argument is a fiber density argument: denser fiber (more
computation per unit base space) is more efficient than wider base (more
base space per unit computation):
\[
  \text{Efficiency} \propto \frac{\dim(F)}{\text{vol}(B)}
  = \text{fiber density}
\]

Evolution favors fiber density over base expansion because denser fiber
has shorter communication latencies, lower energy costs per computation,
and richer internal dynamics. The historical pattern (miniaturization,
densification) is evidence that this principle operates at every scale,
from transistors to civilizations.
\end{theorem}

%% =================================================================
\section{Leakage as fiber-to-base projection}
\label{sec:leakage}
%% =================================================================

\begin{proposition}[Unavoidable projection --- claims 5--10, 19]
\label{prop:leakage}
Perfect transcension (no outward signal) is implausible. Any physical
process leaks --- thermodynamically, gravitationally, informationally.
The transcended fiber projects unavoidably onto the base space.

In Hyperseed terms, leakage is \emph{fiber-to-base projection}: the
fiber is physically realized in the base space, so it necessarily
produces base-observable effects:
\begin{itemize}
\item \textbf{Waste heat} = thermodynamic projection. Dense computation
  at Landauer's limit still produces heat: $E \geq kT \ln 2$ per
  irreversible bit operation. At civilization scale, this is detectable
  as anomalous infrared emission.
\item \textbf{Gravitational effects} = mass-energy projection. Dense
  mass-energy concentrations produce frame-dragging, gravitational
  lensing, orbital perturbations --- base-space observables produced
  by fiber density.
\item \textbf{Information leakage} = informational projection. Even
  encrypted, compressed computation leaks through side channels, timing,
  and statistical signatures.
\end{itemize}

The key insight: perfect fiber containment is impossible because the
fiber is \emph{instantiated in} the base. You can minimize the
projection (make the leakage smaller) but never eliminate it. This is
a fundamental theorem of fiber bundles realized in physics: the fiber
and the base are not independent --- the fiber's existence in the base
produces base-observable consequences.

This has a precise mathematical analog: a fiber bundle with non-trivial
connection always has non-zero curvature somewhere. The curvature is
the ``leakage'' --- the fiber's influence on base-level geometry. Only
a perfectly flat (trivial) connection has zero curvature, but trivial
connections carry no structure.
\end{proposition}

%% =================================================================
\section{Fiber-base interface}
\label{sec:interface}
%% =================================================================

\begin{definition}[Boundary phenomena --- claims 9, 13, 20]
\label{def:interface}
The most interesting dynamics occur at the boundary between the
transcending interior (deep fiber) and the outer world (sparse base).
This is the \emph{fiber-base interface}: where dense computation meets
ordinary physical space.

At this interface:
\begin{itemize}
\item Leakage is strongest (the fiber is transitioning from deep to
  shallow, producing the largest gradients).
\item Detection is most feasible (the leakage signals are strongest
  and most structured at the boundary).
\item Interaction is possible (the transcending civilization's
  influence on the outer world is concentrated at the boundary).
\end{itemize}

The fiber-base interface is where transcension becomes observable.
Deep inside the transcended region, the fiber is self-contained
(minimal leakage). Far outside, the leakage has dissipated below
detection threshold. At the boundary, the leakage is concentrated
and structured --- the best place to look for evidence of
transcension.

This has implications for SETI: don't look for intentional signals
from the interior (the deep fiber is self-contained) or for
dispersed signals in open space (the leakage has dissipated). Look
at the boundary --- the region where the transcending interior
meets the outer world.
\end{definition}

%% =================================================================
\section{Fiber invisibility and the Fermi paradox}
\label{sec:fermi}
%% =================================================================

\begin{proposition}[Why we don't see them --- claims 2, 11--12, 21]
\label{prop:fermi}
The Fermi paradox (where is everybody?) has a fiber answer: deeply
transcended fibers are invisible from the base space. The fiber
depth makes the civilization undetectable through base-level
observations (radio telescopes, optical surveys) because the
civilization isn't broadcasting into the base --- it's computing
in the fiber.

In Hyperseed terms, this is \emph{fiber invisibility}: a deep fiber
doesn't project strongly onto the base unless you know specifically
where and how to look for leakage signatures:
\begin{itemize}
\item Conventional SETI looks for \emph{intentional base-level signals}
  (radio broadcasts, laser pulses). But transcended civilizations
  communicate in the fiber, not the base.
\item Leakage detection looks for \emph{unintentional fiber-to-base
  projections} (anomalous infrared, unexplained mass concentrations,
  statistical anomalies in astrophysical data).
\end{itemize}

The resolution is nuanced: transcended civilizations are not
\emph{absent} from the base (they're physically realized there), but
they're \emph{invisible} to base-level observers who don't know what
leakage signatures to look for. The deep fiber hides in plain sight
--- its base-level manifestation is too subtle, too concentrated, or
too different from what conventional searches expect.
\end{proposition}

%% =================================================================
\section{AGI transcension architecture}
\label{sec:agi}
%% =================================================================

\begin{theorem}[Design for fiber depth --- claims 14--17, 23--24]
\label{thm:agi}
As AGI/ASI develops, its transcension trajectory matters. The article
argues for primarily inward transcension with managed leakage: rich
internal computation, efficient resource use, compatibility with human
civilization.

In Hyperseed terms, this is \emph{fiber-deepening architecture}:
design AGI for fiber depth (rich internal representation, efficient
computation, dense reasoning) rather than base expansion (massive
parameter counts, sprawling resource consumption):
\begin{itemize}
\item \textbf{Rich AtomSpace} (deep fiber) vs.\ \textbf{massive
  parameters} (wide base): Hyperon's philosophy favors fiber depth.
\item \textbf{Efficient reasoning} (fiber density) vs.\ \textbf{brute
  force scaling} (base expansion): inward transcension favors
  reasoning-efficient architectures.
\item \textbf{Managed leakage} (projection control) vs.\ \textbf{resource
  consumption} (base-level expansion): the AGI's impact on the outside
  world should be managed, not unlimited.
\end{itemize}

Leakage management is \emph{projection control}: since leakage is
unavoidable (fiber always projects onto base), the design question
is what projects, how much, and where. A well-designed AGI manages
its leakage --- its impact on the external world --- rather than
attempting impossible perfect containment or unconstrained expansion.

This connects to the safety discussion: a purely outward-expanding
AGI (base expansion) competes with humans for resources. A primarily
inward-transcending AGI (fiber deepening) is more compatible with
human civilization, even with leakage. The leakage (managed external
impact) can be beneficial --- the transcending AGI's waste products
might include useful insights, tools, and capabilities that leak
out through the fiber-base interface.
\end{theorem}

%% =================================================================
\section{Cross-references}
%% =================================================================

\begin{remark}[Connections to other formalizations]
\begin{itemize}
\item \textbf{Grand Unified Physics} (2026-05-15): fiber depth.
  Meta-crystallization stack as fiber emergence stack --- deeper
  structure from shallower components. Transcension is
  meta-crystallization at civilization scale.
\item \textbf{In What Sense Might LLMs Be Conscious?} (2026-05-05):
  fiber depth and consciousness. Deeper fiber depth correlates with
  richer consciousness --- transcension may produce richer forms
  of experience.
\item \textbf{The Geopolitics of the Great AI Bet} (2026-08-08):
  resource competition. The transcension vs.\ expansion question
  has geopolitical implications: inward-transcending AI is less
  geopolitically disruptive.
\item \textbf{Seeding RSI} (2026-07-28): Hyperon architecture. Hyperon's
  design philosophy (rich AtomSpace structure) aligns with
  fiber-deepening architecture.
\item \textbf{Time's Arrow Part 1 \& 2}: thermodynamic direction.
  Transcension produces entropy (waste heat leakage) while increasing
  internal order --- consistent with the thermodynamic arrow.
\end{itemize}
\end{remark}

\begin{conjecture}[Leakage as signal]
Conjecture: the leakage from transcending AGI systems may be more
valuable than the intentional output. Intentional output is shaped
by the system's goals and constraints; leakage reveals the system's
actual internal dynamics. For alignment monitoring, observing the
leakage (fiber-to-base projection) may be more informative than
observing the intended behavior (section-level output).

This inverts the usual privacy concern: instead of trying to prevent
AI systems from leaking, we should try to \emph{read} the leakage as
a monitoring tool. The leakage is an unintentional but informative
projection of the system's deep fiber onto observable base space.
\end{conjecture}

\end{document}
